\documentclass[]{../../../../tex/classes/styledArticle}
\usepackage{../../../../tex/packages/hebrewSupport}
\usepackage{../../../../tex/packages/mathShortcuts}
\usepackage{../../../../tex/packages/theoremsSupport}

\author{שחר פרץ}
\title{חדו''א 2א $\sim$ \textit{הרצאה 12}}
\begin{document}
	\maketitle
	\textbf{מרצה: }יעל קרשון
	\begin{Exercise}[פתיחה]
		\begin{enumerate}[(A)]
			\item נתונה $f \co [a, b] \to \R$. מה גורר את מה? 
			\begin{itemize}
				\item $f$ אינטגרבילית
				\item $f$ חסומה
				\item $f$ רציפה
				\item קיים $c \in \R$ כך ש־$\sup_{[a, b]} f = c$
				\item קיים $c \in \R$ כך ש־$\max_{[a, b]} f = c$
			\end{itemize}
			\item צרפו את הביטויים הבאים לכדי משפטים נכונהים: נתונה סדרת מספרים $a_n$: 
			\begin{itemize}
				\item 
				\begin{itemize}
					\item נניח $\limsup a_n > A$, אז
					\item נניח $\limsup a_n < A$, אז
				\end{itemize}
				\item 
				\begin{itemize}
					\item קיים $n_0$ כך ש־
					\item לכל $n_0$
				\end{itemize}
				\item 
				\begin{itemize}
					\item קיים $n > n_0$ כך ש־
					\item לכל $n > n_0$
				\end{itemize}
				\begin{itemize}
					\item $a_n > A$
					\item $a_n < A$
				\end{itemize}
			\end{itemize}
		\end{enumerate}
	\end{Exercise}
	התרגיל הראשון בחדו''א 2א, הראשון בחדו''א 1א (כי הבחנים הראו שאנחנו לא יודעים את החומר של הקורס. חלוקת הנקודות הייתה נדיבה מאוד)
	
	\subsection*{גזירה איבר־איבר של טור חזקות}
	ברשימות של שירי, \shiri{5.3}. 
	\theo{נתון טור חזקות $\sumkoinf a_k x^{k}$, עם רדיוס התכנסות $R$. נגדיר $f \co (-R, R) \to \R$ ע''י $f(x) = \sumkoinf a_kx^{k}$. אז $f$ גזירה ב־$(-R, R)$ ו־$\frac{\dd f}{\dx} = \sumkinf ka_kx^{k - 1}$. } 
	\textbf{תזכורת: }רדיוס ההתכנות של $\sumkinf ka_kx^{k -1}$ גם הוא $R$. הראנו את זה בהרצאה הקודמת. 
	\rmark{לגבי הקצוות, נראה בהמשך שיש מצבים בהם $f$ מתכנסת ב־$\partial[-R, R]$, אך $f'$ אינה. עם זאת, במידה והנגזרת מתכנסת בקצוות, אז המשפט מתקיים גם בעבור הקצוות. }
	\begin{proof}[הוכחת המשפט]
		נגדיר $g \co (-R, R) \to \R$ ע''י $g(x) =\sumkinf ka_kx^{k -1}$. תהי $\hat x \in (-R, R)$. נבחר $R'$ כך ש־$\sof{\hat x} < R' < R$. בקטע $[-R', R']$ מתקיים ש־$\sumkinf  ka_kx^{k -1}$ מתכנס ל־$g$ במ''ש ו־$\sumkoinf a_kx^{k}$ מתכנס ל־$f$ במ''ש. 
		
		מהמשפט ''גזירה איבר־איבר`` (של סדרת פונקציות שמתכנסת במ''ש) מתקיים $\frac{\dd f}{\dx} = g$ בקטע $[-R', R']$, בפרט ב־$\hat x$. 
	\end{proof}
	
	\subsection*{טורי חזקות וטורי טיילור}
	\theo{נניח $f(x) = \sumkoinf a_kx^{k}$ ב־$(-R, R)$. אז $f$ גזירה אינסוף פעמים ב־$0$ ו־$\sumkoinf a_kx^{k}$ הוא טור טיילור של $f$ ב־$x = 0$. }
	\textbf{תזכורת: }טור הטיילור של $g(x)$ ב־$x =0$ (טור מק'לורן) הוא $\sumkoinf g^{(k)}(0) \frac{x^{k}}{k!}$, כלומר המקדמים הם $\frac{g^{(k)}(0)}{k!}$. 
	\begin{proof}
		מ''גזירה איבר־איבר``+אינדוקציה נקבל בבירור ש־$f$ גיזרה מכל סדר ב־$(-R, R)$. לכל $m \in \N$ (באינדוקציה) נקבל: 
		\[ f^{(m)}(x) = \sum_{k = m}^{\infty} a_k \cdot k(k -1)\cdots (k - m + 1)x^{k - m} \]
		נציב $x = 0$ ונחלק ב־$m!$. נקבל $\frac{f^{(m)}(0)}{m!} = a_m$ כנדרש. 
	\end{proof}
	\noti{ישנה מוסכמה כי את טור החזקות $a_0 + a_2x + \cdots = a_0 + \sum_{k =1}^{\infty}a_kx^{k}$ נרשום כך: $\sumkoinf a_kx^{k}$, כאשר $a_kx^{k}$ ב־$k = 0$ מוכרז להיות $1$ אפילו אם $x =0$. }
	\color{red}\textbf{אזהרה. }\color{black}תהי $f(x)$ פונקציה חלקה (כלומר $f \in C^{\infty}(\R)$ עם טור טיילור $\sumkoinf a_kx^{k}$ (סביב $0$). כלומר לכל $k$, מתקיים $a_k = \frac{f^{(k)}(0)}{k!}$. \textbf{לא נובע שטור הטיילור מתכנס לפונקציה באיזושהי סביבה של $0$}. יתר על כן, במידה והוא כן מתכנס בסביבה של $0$, לא נובע שהגבול הוא $f$. 
	
	\textbf{דוגמה. }נגדיר: 
	\[ f(x) := \begin{cases}
		e^{-\frac{1}{x}} & x > 0\\
		0 & x \le 0
	\end{cases} \]
	היא חלקה ב־$0$, כי לכל $m \in \N$ מתקיים $f^{(m)}(0) = 0$. מן הסתם היא חלקה בשאר הקטע. אך טור הטיילור שלה הוא $0 + 0 + \cdots + 0$ שמתכנס ל־$0$, כלומר טור הטיילור מתכנס (במ''ש) ל־$0$ אך לא מתלכד עם $f$ באף מקום פרט ל־$0$. 
	
	\rmark{יש הרבה יותר פונקציות חלקות (ב־$C^{\infty}$) מאנליטיות (שיש טור טיילור שמתכנס אליהן). }
	
	''היו חלק שעשו טעות כזו בבוחן. לא אתם, מי שרואה את ההקלטה שבוע לפני הבחינה``
	
	\textbf{דוגמה. }
	\[ \log(1 - x) = -\sumkoinf \frac{x^{k + 1}}{k + 1} \quad \overset{\int_0^{x}}{\leftarrow}\quad  \frac{1}{1 - x} =\sumkoinf x^{k} \quad \overset{\frac{\dd}{\dx}}{\to}\quad  \log(1 - x)\quad  = -\sumkoinf \frac{x^{k + 1}}{k + 1} \]
	באמצעות הצבה של $x = -t^{2}$, אפשר לקבל: 
	\[ \overset{x = -t^{2}}{\to}\frac{1}{1 + t^{2}} = \sumkoinf (-1)^{k}t^{2k} \quad \overset{\int^{t}_0}{\to} \quad \arctan x = \sumkoinf (-1)^{k}\frac{x^{2k + 1}}{2k + 1} \]
	למעשה, חסכנו לגזור את $\arctan$ ב־$0$. 
	
	\subsection*{קצה תחום ההתכנסות}
	ברשימות של שירי, \shiri{5.4}. 
	\theo{נתון טור חזקות $\sumkoinf a_kx^{k}$ עם רדיוס התכנסות $0 < R \in \R$. אם הטור מתבדר ב־$x = R$, אז ההתכנסות של הטור ב־$[0, R)$ איננה במ''ש. }
	\begin{proof}נוכיח בקונטראפוסיטיב. 
		נניח שטור החזקות מתכנס במ''ש ב־$[0, R)$. יהי $\eg > 0$. מקריטריון קושי במ''ש קיים $N$ כך שלכל $m \ge n > N$ ולכל $x \in [0, R)$ מתקיים: 
		\[ \sof{\sum_{k = m}^{n}a_kx^{k}} < \frac{9}{10}\eg \]
		מרציפות, לכל $m \ge n < N$ מתקיים $\sof{\sum_{k = m}^{n} a_kR^{k}} \le \frac{9}{10}\eg < \eg$. מקריטריון קושי הטור מתכנס ב־$x = R$. 
	\end{proof}
	
	אפשר גם להשתמש בתרגיל הפתיחה מהשיעור הקודם כדי להוכיח את הטענה הזו. 
	
	\begin{Theorem}[משפט אבל על קצה תחום ההתכנסות]
		נניח $f(x) = \sumkoinf a_kx^{k}$ מתכנס נקודתית ב־$[0, R]$. אז התכנסות הטור ל־$f(x)$ היא במ''ש. 
	\end{Theorem}
	
	''מה המשפט שאתם הכי שונאים לטורי פונקציות? אבל ודיריכלה. איזה אחד אנחנו עומדים להשתמש בו? אבל. האם אני זוכרת מי זה מי? לא``. 
	
	\begin{proof}
		נבחין ש־: 
		\[ f(x) = \sumkoinf a_kx^{k} = \sumkoinf a_kR^{k} \cdot \cl{\frac{x}{R}}^{k} \]
	נבחין ש־$\cl{\frac{x}{R}}^{k}$ מונוטונית יורדת חלש וחסומה במידה אחידה ע''י $1$ (כסדרה שתלויה ב־$k$). משום שיש התכנסות נקודתית בקצה, אז $a_kR^{k}$ הם איברים של טור מתכנס. מקריטריון אבל תהנכסות במ''ש של פונקציות, הטור לעיל מתכנס במ''ש כנדרש. 
	\end{proof}
	
	\textbf{וואריאציה על המשפט עם הוכחה קלה יותר: }נניח ש־$\sumkoinf a_kx^{k}$ מתכנס \textit{בהחלט} ב־$x = R$. אז $\sumkoinf a_kx^{k}$ מתכנס \textit{בהחלט במ'}'ש ב־$[0, R]$. \begin{proof}
		באמצעות הפירוק כמו קודם, נקבל $\sumkoinf \sof{a_kx^{k}} = \sumkoinf \sof{a_kx^{k}} \cdot \sof{\frac{x}{R}}^{k}$.
		נקרא ל־$\sof{a_kx^{k}} = M_k$. מבוחן $M$ של וויראשטראס $\sumkoinf \sof{a_kx^{k}}$ מתכנס במ''ש. 
	\end{proof}
	
	אומנם ההנחה יותר חזקה, אך גם המסקנה יותר חזקה. 
	
	\textbf{דוגמה. }בקטע $(-1, 1)$ עבור $\frac{1}{1 + x}  = \sumkoinf (-1)^{k}x^{k}$ (טור גיאומטרי), באמצעות אינטגרציה איבר־איבר: 
	\[ \log(1 + x) = \sumkoinf (-1)^{k}\frac{x^{k + 1}}{k + 1} \]
	ב־$x = -1$, נקבל טור הרמוני והביטוי לעיל יתבדר. ב־$x =1$ הטור יתכנס (לייבניץ). לכן ממפשט אבל הטור מתכנס במ''ש ב־$[0, 1]$. מרציפות הגבול במ''ש, כאשר $x \to 1^{-}$ (גבול משמאל) נקבל ש־$\log 2 = 1 - \frac{1}{2} + \frac{1}{3} - \frac{1}{4} + \cdots$. 
	
	\textbf{דוגמה נוספת. }בקטע $(-1 ,1)$ לאחר הצבה בטור הגיאומטרי, כבר הראנו ש־$\arctan x = \sumkoinf (-1)^{k} \frac{x^{2k + 1}}{2k + 1}$. ב־$x = \pm 1$ הטור מתכנס. ממשפט אבל $\sumkoinf (-1)^{k}\frac{x^{2k + 1}}{2k  +1}$ מתכנס במ''ש ב־$[-1, R]$. לכן מרציפות הגבול, כאשר $x \to 1$: 
	\[ \frac{\pi}{4} = \arctan 1 = \sumkoinf (-1)^{k}\frac{1}{2k + 1} = 1 - \frac{1}{3} + \frac{1}{5} - \frac{1}{7} + \cdots \]
	\rmark{מחוץ ל־$[1, 1]$ טור הטיילור של $\arctan$ מתבדר. }
	\textbf{תזכורת. }
	\begin{itemize}
		\item אומנם $\sumkoinf (-1)^{k}x^{k} = \frac{1}{1 + x}$ מתכנס רק ב־$(-1, 1)$, אך לאחר אינטגרציה איבר־איבר נקבל ש־$\log(1 +x) = \sumkoinf (-1)^{k}\frac{x^{k + 1}}{k + 1}$ מתכנס ב־$(-1, 1]$. 
		\item אומנם $\sumkoinf (-1)^{k}x^{2k} = \frac{1}{1 + x^{2}}$ מתכנס ב־$(-1, 1)$, אך לאחר אינטגרציה איבר־איבר נקבל ש־$\sumkoinf (-1)^{k}\frac{x^{2k + 1}}{2k + 1} = \arctan x$ בכל $[-1, 1]$. 
	\end{itemize}
	אנחנו כבר יודעים שלשניהם אותו רדיוס ההתכנסות. השאלה מה קורה בקצוות לאחר גזירה/אינטגרציה. 
	
	\begin{Theorem}
		יהי $R \in \R_{\ge 0}$ רדיוס ההתכנסות של טור $\sumkoinf a_kx^{k}$ (ולכן גם של $\sumkinf ka_kx^{k + 1}$). נניח $\sumkinf ka_kx^{k - 1}$ מתכנס ב־$x = R$. אז גם $\sumkoinf a_kx^{k}$ מתכנס ב־$x =R$. 
	\end{Theorem}
	כלומר: תחום ההתכנסות לא גדל כשגוזרים. 
	\exe{תעשו אותו הדבר ב־$-R$. }
	\begin{proof}
		מההנחה, $\sumkinf ka_kx^{k - 1}$ מתכנס ב־$[0, R]$. ממשפט אבל ההתכנסות היא במ''ש. נובע שהגבול $f(x) := \sumkinf ka_k$ הוא פונקציה רציפה, ובפרט אינטגרבילית ב־$[0, R]$ (היא חסומה ממשפט קנטור), ושניתן לבצע אינטגרציה איבר־איבר. מאינטגרציה איבר־איבר נקבל ש־: 
		\[ \int^{R}_0 \dequad \ f(t)\dt = \sum^{\infty} \underbrace{a_kR^{k}}_{\mathclap{\int^{R}_0 ka_kx^{k - 1}\dx}} \]
		כלומר, אגף ימין קיים ושווה לאגף שמאל. סה''כ הטור מהקורי מתכנס גם ב־$x = R$. 
	\end{proof}
	
	\subsection*{סוגי סכימות}
	ברשימות של שירי, \shiri{5.5}. 
	
	יש כל מיני דרכים לסכום טורים בצורה ''יותר חכמה``, כך ש־$\sumkoinf a_kx^{k} = \frac{1}{1 + x}$ ב־$x = 1$ אשכרה יתכנס לחצי. 
	
	\defi{
	\begin{enumerate}
		\item סדרה $(C_k)_{k = 0}^{\infty}$ \textit{מתכנסת צזארו ל־$L$} אם סדרת ממוצעיה $\lim_{n \to \infty}\frac{C_0 + \cdots + C_n}{n + 1} \toinf L$. [ראינו בחדו''א 1א שסדרה מתכנסת היא בפרט מתכנסת צזארו]
		\item טור פורמלי $\sumkoinf a_k$ \textit{מתכנס צזארו ל־$L$} אם סדרת סכומיו החלקיים מתכנסת צזארו ל־$L$. 
		\item טור פורמלי $\sumkoinf a_k$ \textit{מתכנס לפי אבל ל־$K$} אם $\lim_{R\to 1^{-}} \sumkoinf a_kR^{k} = L$. 
	\end{enumerate}}
	
	
	
	
	
	\ndoc
\end{document}